

\chapter{Introdução}
\label{chapter:introduction}

\begin{introduction}
O projeto nasceu de um problema concreto: a ausência de equipamento de ensaio sísmico acessível em países com recursos limitados. Este capítulo apresenta esse contexto, define os objetivos que guiaram o desenvolvimento e descreve a organização do documento.
\end{introduction}

\section{Motivação e Contexto}
\label{sec:motivacao}

Portugal situa-se numa zona de atividade sísmica moderada a elevada. A sua posição geográfica, próxima da fronteira entre as placas Euroasiática e Africana, cujo limite ocidental é a falha Açores-Gibraltar~\cite{Borges2001}, explica parte substancial dessa perigosidade. Foi essa mesma configuração que gerou o sismo de 1 de novembro de 1755, estimado com magnitude entre 8,5 e 8,7, que destruiu grande parte de Lisboa e desencadeou um tsunami com consequências devastadoras~\cite{Baptista1998}.

O problema não é exclusivo de Portugal. Muitos países, sobretudo os menos desenvolvidos, enfrentam riscos sísmicos comparáveis sem dispor de meios para preparar engenheiros capazes de os avaliar e mitigar. O trabalho de avaliação de risco estrutural no terreno exige acesso a equipamento de ensaio sísmico experimental, mas o custo das soluções comerciais disponíveis coloca-as fora do alcance desses contextos.

Foi precisamente essa lacuna que motivou o presente projeto. O Professor Vítor Silva deparou-se, durante atividades de investigação e cooperação técnica no estrangeiro, com a ausência de equipamento acessível para realizar ensaios sísmicos portáteis de apoio à avaliação de risco estrutural em países com recursos limitados. As mesas sísmicas disponíveis no mercado, como a \textit{Quanser Shake Table II XY}~\cite{QuanserShakeTableIIXY}, custam mais de dez mil euros. O objetivo do presente projeto é, por isso, desenvolver uma alternativa de baixo custo, reprodutível e de arquitetura aberta.

\section{Objetivos do Projeto}
\label{sec:objetivos}

O objetivo central é desenvolver uma mesa sísmica de dois eixos, horizontal~(X) e transversal~(Y), de baixo custo, recorrendo predominantemente a peças fabricadas por impressão 3D e a componentes eletrónicos de uso geral. O sistema tem de ser funcional, reprodutível e suficientemente acessível para uso em investigação e avaliação de risco sísmico com recursos limitados. Os objetivos específicos são:

\begin{itemize}
    \item Projetar e fabricar, por impressão 3D, os componentes mecânicos estruturais da mesa sísmica, com rigidez suficiente para suportar o movimento bidirecional controlado;
    \item Implementar um sistema de atuação baseado em quatro motores de passo NEMA17, controlados por \textit{drivers} DRV8825, com transmissão de movimento por fuso trapezoidal de \textit{lead} de 8\,mm por volta;
    \item Desenvolver o firmware para o microcontrolador ESP32, com recurso à biblioteca AccelStepper, com uma janela de execução de 20\,ms por ponto de dados sísmicos e capacidade para reproduzir registos sísmicos;
    \item Assegurar a reprodutibilidade do sistema, pela disponibilização pública de todos os ficheiros necessários à construção de uma unidade idêntica, com custo total significativamente abaixo das soluções comerciais equivalentes.
\end{itemize}

